%% kapa: exam=1 round=35 year=2024 period=1 form=A subject=civil_law subject_ko=민법 seq=29 number=29
%% source: Q-Net 감정평가사 기출문제 (cst00309) artlSeq=5213229 file=raw/1008/5213229/1교시 문제지 원본.pdf
%% conversion: kordoc 4.24.51 -> tools/gen1.py (auto); review=proofread-sonnet; flags=none
\begin{question}{29}
물권의 소멸에 관한 설명으로 옳지 않은 것은? (다툼이 있으면 판례에 따름)
\begin{choices}
  \item X토지에 甲이 1번 저당권, 乙이 2번 저당권을 취득하고, 丙이 X토지를 가압류한 후 乙이 X토지를 매수하여 소유권을 취득한 경우 乙의 저당권은 혼동으로 소멸하지 않는다.
  \item 유치권자가 유치권 성립 후에 이를 포기하는 의사표시를 한 경우에도 점유를 반환하여야 유치권은 소멸한다.
  \item 점유권과 소유권은 혼동으로 소멸하지 아니한다.
  \item 지역권은 20년간 행사하지 않으면 시효로 소멸한다.
  \item 후순위 저당권이 존재하는 주택을 대항력을 갖춘 임차인이 경매절차에서 매수한 경우, 임차권은 혼동으로 소멸한다.
\end{choices}
\end{question}
\answer{2}
